\chapter{Conclusions and Future Work}
\label{ch:conclusions}

\section{Summary}
\label{sec:summary}

The goal of this project was to design, implement and evaluate a full-duplex
speech-to-speech prototype for Persian. Two tracks were pursued in parallel: adapting an
end-to-end speech-to-speech model, and a modular cascade with full-duplex control. The
outcome was asymmetric.

The \textbf{streaming cascade} reached a working prototype that operates with an open
microphone: it detects the end of speech automatically and speculatively, generates and
plays the reply as a stream, clause by clause, and reacts to the user's interruption during
playback. In loopback testing with a real browser, the mean latency from end of speech to
the assistant's first audio was 372 ms and its 90th percentile 431 ms, and none of the 201
turns went unanswered or exceeded 500 ms; the GPU side of the whole system used 11.1~GiB of
memory. The 40-row final semantic test of the Qwen3-4B responder passed all seven automatic
conditions (relevance 3.150 against 1.425 for the baseline).

\textbf{Interruption detection} reached the accuracy goal with a two-stage controller: one frame of
speech louder than the expected echo ducks the assistant's voice, and then the speech recognizer,
after setting aside listener tokens, function words and the assistant's own echo, decides to stop
or continue. In a deployment simulation built from real recordings, an accuracy of 0.933 and an F1
of 0.880 with a median of 120 ms to the duck were obtained, and in a real browser an accuracy of
0.825 and an F1 of 0.817. The energy/pitch/MFCC acoustic classifier alone was limited to 0.574 on
real podcast events; this negative result is the reason for the lexical design.

The \textbf{components} were re-selected on fresh panels under pre-registered protocols: the NeMo
recognizer was the most accurate candidate within the latency budget on human transcripts, the
Mana voice the most intelligible voice to two automatic listeners (against the other Piper voices,
ParsVoice-XTTS and the lab VITS model, even after optimizing its inference), and no replacement
responder reached the replacement threshold.

The \textbf{direct track} did not reach a deployable model, but its main cause was found: Moshi's
English tokenizer needs 6.6 tokens per Persian word, and 98.5\% of turns exceed the
one-token-per-frame budget. Removing this bottleneck with phoneme romanization doubled the best
generation result (6/9 against 3/9) and turned silent output into audible speech, but the speech
was still not intelligible. The experience of J-Moshi~\cite{ohashi2025jmoshi} confirms this
diagnosis: bringing Moshi to Japanese also began with replacing the tokenizer, but then needed about
60,000 hours of Japanese dialogue pre-training, more than five hundred times the 108.6 hours of this
project's data. The speech-input Gemma~4 model, even with a self-distillation adapter, was weaker
than the cascade.

The \textbf{dataset} is 6,754 ``user speech and response'' pairs with 108.584 hours of audited audio
and zero group leakage. Its labels are automatic and have no human labels.

\section{Answer to the Research Question}
\label{sec:answer}

The practical question of this project had two parts. Part one: with the budget of a bachelor's
project, how far can one go in building a full-duplex Persian speech-to-speech system? The answer
is that a streaming cascade with interruption control meets the proposal's numeric criteria for
latency (mean and 90th percentile below 500 ms) and interruption detection (above 80\%) in automated
tests and fits in the memory of a 12 GB card; but an end-to-end Persian model was not achievable with
the existing architectures and tokenizers within this budget.

Part two: what can be proven with auditable evidence? The answer is in
Table~\ref{tab:req-status}. R1 and R4 are met; R2, R3 and R5 are met with automatic evidence
(browser loopback and simulation), R3 by the proposal's criterion and, for the main panel, even by
the every-turn reading, and R5 by a method other than the specified acoustic features; R6 is a
declared deviation (memory was measured, speed on target hardware was not); and R7 depends on the
responses of the listening-study participants. The methodological contribution of this report lies
in this separation: for each measure the permitted claim boundary is stated, and negative results are
reported as fully as positive ones.

\section{Answers to the Research Questions of the Proposal}
\label{sec:rq-answers}

\textbf{Question 1: low-cost adaptation of an existing speech-to-speech model.} Low-rank adaptation of
Moshika with a LoRA adapter on an H100 improved the training objective: on the held-out test, the
perplexity of the text tokens went from 2.53 to 1.19 and that of the audio tokens from 13.98 to 4.71.
But in none of the seven generations was intelligible Persian speech produced. The main cause is the
tokenizer: the English tokenizer needs 6.6 tokens per Persian word, while the model produces at most
one text token per 12.5 Hz frame. Phoneme romanization brought this load to 0.67 tokens per frame and
doubled the best run (6/9), but it was not enough. Answer: not with ``the least retraining cost'';
adapting this architecture to Persian needs at least a text representation compatible with Persian
and training beyond a small adapter.

\textbf{Question 2: architectural changes needed for full duplex below 500 ms.} In the cascade, these
changes were necessary and sufficient: (1) speculative end-of-speech detection that starts the work
before the end is confirmed and holds the audio until confirmation; (2) streaming token generation
with a ``Persian script only'' constraint at decoding time; (3) clause chunking that builds the first
fragment from two to four words and synthesizes each clause immediately; (4) a playback queue in the
client and an audio-holding gate in the server, which together play the role of a ``dual buffer'';
and (5) two-stage interruption control: ducking within one frame and lexical confirmation that sets
aside listener tokens and echo, together with keeping only the heard part of an interrupted reply in
the history. With this design, the mean latency was 372 and the 90th percentile 431 ms, and the reply
after an interruption stayed as relevant as the reply to a normal turn
(Section~\ref{sec:res-latency}). In the direct track, the change needed before any architectural
change is the tokenizer.

\textbf{Question 3: the trade-off between latency, size and quality.} In this project size did not buy
quality (Table~\ref{tab:tradeoff}): the smallest recognizer fine-tuned on Persian (117 million
parameters) was more accurate than all the general models of one to nine billion parameters, and the
only model more accurate than it, itself fine-tuned on Persian, was more than ten times slower; the
smallest synthesizer (15.7 million parameters) was both the most intelligible and the fastest; and the
largest model, the 7.7-billion-parameter end-to-end model, did not reach intelligible speech. The best
trade-off point for real-time Persian is small components specialized for Persian.

\textbf{Question 4: the perception of older adults.} The instrument and protocol of the listening study
are ready and pre-registered, but the participants' responses have not yet arrived; until then this
report makes no claim about the perception of older adults.

\section{Future Work}
\label{sec:future}

\subsection{Completing the Evidence}

\begin{enumerate}
\item \textbf{Target hardware and a physical path.} Running the same loopback protocol on a 12--24 GB
      card, and a test with a physical microphone and loudspeaker that includes real echo and the
      browser's echo cancellation.
\item \textbf{Human study.} Collecting the responses of 5--10 participants (at least two aged over 60)
      with the registered instrument; and as a next step, a short live interaction session with older
      adults.
\item \textbf{Human labels.} A listening review of a sample of interruption events and dataset pairs by
      two independent reviewers, to validate the automatic Gemini labels and the extraction rules.
\item \textbf{Human judging of the responder.} The two automatic judges disagree considerably about the
      responders (Gemini rates the responders with the elderly prompt about one point better and the
      Qwen judge does not). The replacement rule was registered before the listening study was sent,
      and the configuration of Qwen3.5-4B with the elderly prompt has passed its deployment
      qualification; the participants' preference therefore decides whether it replaces the deployed
      responder.
\end{enumerate}

\subsection{Improving the System}

\begin{enumerate}
\item \textbf{More accurate lexical confirmation.} The first phoneme of short backchannels is sometimes
      clipped (\fa{«آره»} as \fa{«کار»}). Adding 200 ms of audio before the onset to the recognizer's
      input was tested and rejected: more interruptions were stopped (0.98 versus 0.95) but more
      backchannels and noise too, and accuracy fell (Section~\ref{sec:res-bargein-sim}). Separating
      listener words from content words with a clipped onset needs labelled data and a learned model.
\item \textbf{End-of-speech detection for older adults.} The long pauses of slow speech sometimes close
      the turn before the end of the sentence (87 server turns for 80 planned turns). The system
      compensates when the user continues by stopping the reply. A semantic end-of-turn gate was tested
      in two rounds (Section~\ref{sec:res-latency}): the best one, an LLM completeness judgement,
      split fewer utterances in the loopback test but held 15\% of finished turns and pushed the 90th
      percentile above 500 ms. A gate that holds for less time, or a Persian end-of-turn model trained
      on turn ends labelled by people rather than on podcast speaker boundaries, is the next step.
\item \textbf{A more accurate recognizer within the same budget.} Seamless had a lower character error
      rate but more than four times the latency budget; distilling it into a CTC model or running it in
      streaming mode could raise accuracy without a latency cost.
\end{enumerate}

\subsection{The Direct Track}

The tokenizer finding points to a clear path: a direct Persian model needs a tokenization that fits
within the one-token-per-frame budget (phoneme-based or with a Persian vocabulary), pre-training on
Persian dialogue at a scale close to J-Moshi's (tens of thousands of hours, for which no public
Persian corpus exists), and a remedy for the gap between teacher forcing and free generation, such as
scheduled input dropout at a larger scale or training on the model's own output. The frozen-model
approach of Freeze-Omni~\cite{wang2025freezeomni} and stabilizing Persian speech generation with
plentiful single-speaker data such as ManaTTS~\cite{qharabagh2025manatts} before dialogue training are
also promising paths. For the speech-input model, training on more varied speech (not only podcasts)
is needed, because self-distillation on one domain hurt outside that domain.

\section{Closing Remarks}
\label{sec:closing}

This project showed that a full-duplex Persian voice assistant with an open microphone, latency below
half a second and reliable interruption detection can be built within the budget of a bachelor's
project with open components; and it showed that an end-to-end Persian speech-to-speech model needs
more than low-rank adaptation, including a tokenization suited to the language. The negative result
of the direct track, together with its measured cause, is worth no less than the positive result of
the cascade: offline loss reduction can be misleading, and a runtime generation test must be part of
the selection criterion.
